\section{Introduction }

The previous chapter defined AI-Scientist as a verification-first, evidence-grounded research assistant. This
chapter shows how that definition is realized architecturally. The most important high-level point is that
AI-Scientist is not a single prompt or an unstructured collection of agents \cite{wu2023autogen,li2023camel,du2023debate}. It is a structured pipeline of
specialized agents coordinated by an explicit Orchestrator, in which control flow is deterministic and only
specific stages delegate to a language model. This design keeps the system focused, inspectable, and well suited
to thesis-level evaluation.

The architecture follows the evidence discussed in the literature review. Grounded reasoning systems are easier to
defend, inspect, and reason about when their stages and transitions are visible. AI-Scientist therefore separates
retrieval, claim extraction, verification, critique, synthesis, and reporting into distinct agents with typed
inputs and outputs. This separation gives the thesis a clear systems story: trustworthy scientific answering is not
a monolithic generation step, but a controlled workflow that transforms a research question into an
evidence-backed, confidence-annotated report.

\section{Architectural Philosophy }

AI-Scientist is best understood as a \emph{multi-agent orchestrated pipeline} with an embedded re-verification
loop \cite{yao2023treeofthoughts,shinn2023reflexion,asai2024selfrag}. A single Orchestrator owns the workflow state and invokes each agent in turn; the agents do not coordinate
autonomously among themselves. In the terminology used throughout the project:

\begin{enumerate}
\item an \textbf{agent} is a component with one specialized responsibility (retrieval, extraction, verification,
critique, synthesis, or reporting);
\item a \textbf{deterministic stage} uses fixed rules, token-overlap measures, and thresholds and does not call a
language model;
\item an \textbf{LLM-backed stage} performs bounded model-assisted refinement on top of a deterministic result;
and
\item the \textbf{critic loop} is the adaptive act-observe-revise step in which the Critic's assessment of weak
claims feeds a second, targeted retrieval-and-verification pass.
\end{enumerate}

This distinction matters because it clarifies where each capability contributes most effectively. The backbone of
AI-Scientist --- retrieval ranking, evidence scoring, verdict selection, confidence computation, and critique ---
is deterministic and therefore reproducible and inspectable. Language-model assistance is layered on selected
stages as an optional refinement, and the system degrades gracefully to its deterministic behaviour when no model
is available. This separation mirrors the direction of agentic and retrieval-augmented systems that favor explicit control surfaces over a single opaque generation step \cite{lewis2020rag,guu2020realm,langgraph2024}.
Figure~\ref{fig:architecture} shows the high-level organization, and Figure~\ref{fig:pipeline} shows the
end-to-end data flow including the critic loop.

The overall division of responsibilities is summarized in Table~\ref{tab:agent_roles}.

\begin{table}[!htbp]
\centering
\caption{Specialized agents and their primary responsibilities}
\label{tab:agent_roles}
\begin{tabularx}{\linewidth}{>{\raggedright\arraybackslash}p{0.25\linewidth} >{\raggedright\arraybackslash}X}
\toprule
Agent & Primary responsibility \\
\midrule
Orchestrator & Coordinates workflow state, scope control, and stage transitions \\
Retrieval & Finds and ranks evidence from uploads, corpus, and live sources \\
Claim Extraction & Converts evidence into checkable scientific claims \\
Verification & Assigns verdicts and confidence using evidence support \\
Critic & Identifies weak claims and triggers targeted re-verification \\
Synthesis & Produces the concise final answer from verified claims \\
Report & Assembles the traceable structured output for the user \\
\bottomrule
\end{tabularx}
\end{table}

\begin{figure}[!htbp]
\centering
\includegraphics[draft=false,width=0.72\linewidth]{Fig_4.1.jpg}
\caption[High-level AI-Scientist architecture]{\textbf{High-level AI-Scientist architecture.} A single
Orchestrator coordinates the specialized agents and produces an evidence-backed report with verdicts,
confidence, and traceability.}
\label{fig:layered_system_view}
\label{fig:architecture}
\end{figure}

\section{The Orchestrator Agent }

The Orchestrator is the central controller of AI-Scientist. It receives the user's research question, applies the
scope guard, and then drives the agents in a fixed order: retrieve evidence, extract claims, verify each claim,
criticize the results, optionally re-verify weak claims, synthesize a summary, and build the report \cite{langgraph2024,wu2023autogen}.
It also owns
the workflow's early-exit conditions: if the question is out of scope, if no relevant papers are found, or if no
sufficiently aligned claims can be extracted, the Orchestrator stops early and returns an honest report explaining
why, rather than forcing the pipeline forward into an uncertain conclusion.

Crucially, the Orchestrator maintains the \emph{orchestration trace} --- a running list of human-readable steps
describing what happened at each stage, such as how many papers were retrieved, how many claims were extracted,
how many critique notes were raised, and whether a second verification pass was triggered. This trace is part of
the final report and is the mechanism by which the system's reasoning becomes inspectable.

This trace also creates a convenient thesis artifact. It turns the internal workflow into a visible sequence that
can be discussed in text, displayed in the interface, and examined during evaluation. In that sense, the
Orchestrator is not just a control component; it is the layer that makes the entire pipeline legible to the user.

\section{The Evidence Layer and the Retrieval Agent }

The Retrieval agent realizes the layered evidence model from Chapter~3. It maintains a local pool that lists
user-uploaded papers first and then the curated corpus, deduplicated so that an uploaded paper wins on a tie. When
live retrieval is enabled, it augments this local pool with results fetched at query time from PubMed Central and
arXiv \cite{lewis2020rag,guu2020realm,karpukhin2020dpr,khattab2020colbert}.

Ranking is where the source-priority requirement becomes concrete. Each candidate paper is scored by its
query-coverage relevance plus a source-dependent bonus, with uploaded papers receiving the largest bonus, curated
corpus papers an intermediate bonus, and live-API papers the smallest. Candidates are then admitted only if their
raw coverage clears the minimum threshold, so the source bonus changes ordering but cannot smuggle in an
irrelevant paper. The agent also applies a quality filter to live results --- penalizing very short or metadata-poor
abstracts --- so that low-quality API hits do not dilute the evidence pool. The result is a ranked evidence set in
which user-provided and locally trusted papers surface first, while genuinely relevant live results can still
appear.

This ranking logic is important because it gives the architecture a practical notion of evidence preference. It
allows the system to stay responsive to the user's own materials, remain grounded in a reliable offline corpus, and
still extend outward when the question benefits from broader coverage. The evidence layer therefore supports both
precision and flexibility.

\section{The Claim Extraction Agent }

Once evidence has been retrieved, the Claim Extraction agent converts the unstructured abstracts into discrete,
checkable claims \cite{thorne2018fever,wadden2020scifact}. It scans each retrieved paper sentence by sentence and selects candidate statements that either
contain research-assertion cues (such as ``improves,'' ``reduces,'' ``demonstrates,'' or ``outperforms'') or share
content tokens with the user's question. Each selected sentence becomes a structured claim that records its text, a
coarse claim type (performance, finding, implication, limitation, or observation), its source paper, and the focus
terms it shares with the question.

This stage is the architectural embodiment of the claim-level factuality principle from the literature. By turning
a paper into a small set of explicit claims, AI-Scientist ensures that verification, confidence, and critique all
operate at the granularity where support and contradiction actually differ, rather than on a single fused answer.
The number of claims per question is bounded so that the workflow stays focused on the statements most relevant to
the query.

This bounded extraction stage has two practical benefits. First, it keeps the downstream verifier focused on the
claims that matter most for the question. Second, it gives the thesis a clean object of analysis: each extracted
statement can be tracked, scored, and discussed individually. That makes the system easier to evaluate and the
writing easier to defend.

\section{The Verification Agent }

The Verification agent is the analytical core of AI-Scientist. For each claim, it examines the sentences of the
candidate papers and computes an evidence score that blends Jaccard overlap with focus-term and content-term
agreement, plus small bonuses when a sentence comes from the claim's own source or matches its rhetorical type. It
then classifies each scored sentence as supporting or contradicting the claim, using negation and
contradiction-cue analysis: a sentence whose polarity disagrees with the claim's polarity is treated as
contradictory.

From the collected evidence, the agent selects one of the four verdicts defined in Chapter~3. A claim with
adequate aligned support becomes supported; a claim where strong, multi-source counter-evidence dominates becomes
contradicted; an overgeneralized claim, or one without sufficient overlap, becomes insufficient evidence. The
confidence attached to each verdict is computed deterministically from evidence strength, the number of additional
supporting snippets, and the number of independent corroborating sources, with an explicit penalty when a claim is
only self-corroborated. This is the architectural point at which the honesty requirements become operational:
verdicts are bound to evidence, contradictions are representable, and self-corroboration is discounted.

The deterministic confidence computation is valuable from a thesis perspective because it produces a stable and
interpretable signal. The score is not arbitrary; it reflects observable evidence properties that can be explained
clearly in text, compared across claims, and discussed during evaluation.

When a language model is available, the verifier can pass its preliminary verdict, together with the selected
supporting and contradicting snippets, to a bounded refinement step that may adjust the verdict and rationale. This
refinement is constrained to the already-retrieved evidence and is optional; if it is unavailable or fails, the
deterministic assessment stands.

\section{The Critic Agent and the Re-Verification Loop }

The Critic agent implements the adaptive-reasoning requirement \cite{shinn2023reflexion,yao2023treeofthoughts,asai2024selfrag}. After the first verification pass, it reviews all
assessments and raises critique notes for four weakness patterns: claims with insufficient evidence, claims whose
confidence falls below a threshold, claims supported by only a single source, and claims that are contradicted.
These notes are reported to the user, but more importantly they can drive action.

For claims that are weak in a way that more evidence might resolve, the Critic produces a re-verification decision:
a targeted follow-up query built from the claim text, its focus terms, and its source title, with an added
counter-evidence probe when the claim contains absolute or contradiction-prone language. The Orchestrator then
runs a second retrieval pass for that specific claim, merges the new papers with the first-pass evidence, and
re-verifies the claim with a preference for resolving contradictions. Overgeneralized claims are deliberately
excluded from this loop, because seeking more evidence cannot rescue a claim that is too broad to settle.

This loop is the architectural difference between AI-Scientist and a fixed RAG pipeline \cite{lewis2020rag,ragsurvey2024}. The effort spent on a
claim is conditioned on how uncertain the first pass was, and the additional retrieval is aimed precisely at the
weak claim rather than re-running the whole query. Figure~\ref{fig:pipeline} shows this feedback path explicitly.

The loop is also thesis-friendly because it demonstrates controlled adaptivity. The system responds to weak
evidence by deciding where additional effort should be spent. That makes the design feel deliberate and
research-oriented.

\begin{figure}[!htbp]
\centering
\begin{tikzpicture}[
  font=\small,
  node distance=5.5mm,
  scale=0.95,
  transform shape,
  stage/.style={draw, rounded corners, align=center, text width=26mm, minimum height=9mm, inner sep=3pt, fill=blue!15},
  io/.style={draw, rounded corners, align=center, text width=26mm, minimum height=9mm, inner sep=3pt, fill=orange!20},
  >={Stealth[round]}
]
\node[io] (q) {Research Question};
\node[stage, below=of q] (scope) {Scope Guard};
\node[stage, below=of scope] (ret) {Retrieval\\ \scriptsize (uploads / corpus / live)};
\node[stage, below=of ret] (claim) {Claim Extraction};
\node[stage, below=of claim] (ver) {Verification\\ \scriptsize (verdict + confidence)};
\node[stage, below=of ver] (crit) {Critic\\ \scriptsize (flag weak claims)};
\node[stage, below=of crit] (syn) {Synthesis};
\node[io, below=of syn] (rep) {Report + Trace};

\draw[->] (q) -- (scope);
\draw[->] (scope) -- (ret) node[midway,right]{\scriptsize in scope};
\draw[->] (ret) -- (claim);
\draw[->] (claim) -- (ver);
\draw[->] (ver) -- (crit);
\draw[->] (crit) -- (syn) node[midway,right]{\scriptsize accepted};
\draw[->] (syn) -- (rep);

% critic re-verification loop
\draw[->, dashed] (crit.east) -- ++(16mm,0) |- (ret.east)
  node[pos=0.25, right, align=left]{\scriptsize targeted\\ \scriptsize re-verification};

% early exits
\draw[->, dotted] (scope.west) -- ++(-12mm,0) |- (rep.west)
  node[pos=0.25, left, align=right]{\scriptsize out of scope /\\ \scriptsize no evidence};
\end{tikzpicture}
\caption[The AI-Scientist verification pipeline]{\textbf{The AI-Scientist verification pipeline end-to-end.} A
research question passes the scope guard, is grounded by layered retrieval, decomposed into claims, and verified
into per-claim verdicts with confidence. The Critic flags weak claims and can trigger a targeted re-verification
pass (dashed). Honest early exits (dotted) return a report when the question is out of scope or no relevant
evidence exists. Synthesis and reporting produce the final evidence-backed report and orchestration trace.}
\label{fig:pipeline}
\end{figure}

\section{The Synthesis Agent }

The Synthesis agent composes the executive summary from the verified claims and the critique notes. Its
deterministic form aggregates the assessments into a short, calibrated statement: how many claims are supported,
how many are contradicted, and how many remain unresolved, with an explicit note that weakly supported or
contradicted claims should be read cautiously. When a language model is available, the agent can rewrite this
summary into more fluent prose under tight constraints --- it is instructed to state the overall finding in a few
sentences, to remain grounded in the verified claims, and not to overclaim --- and it falls back to the
deterministic summary if the model is unavailable.

The architectural significance of this stage is that synthesis is downstream of verification, not a substitute for
it. The summary describes the verified evidence and does not introduce new conclusions that were not checked.

This stage directly supports the user-facing answer that appears in the interface. It bridges the gap between
claim-by-claim verification and a concise conclusion, so the final report remains readable without losing its
evidence trail. In thesis terms, it shows how a structured system can be both rigorous and approachable.

\section{The Report Agent and the Orchestration Trace }

The Report agent assembles the final structured output: the question, the executive summary, the list of retrieved
papers, the orchestration trace, the per-claim verification details (claim text, type, source, verdict, confidence,
rationale, and evidence snippets), and the critique notes. The report is produced as both a structured object and a
rendered Markdown document so that it can be displayed in the user interface, returned through the API, or exported.

The trace-based design is intentional. Because the thesis emphasizes trustworthy reasoning, the system should
output more than a conclusion. The conclusion must be inspectable, explainable, and attributable to its
intermediate stages. A reader can see why a claim was marked insufficient, which paper supported a supported claim,
and whether a second verification pass was run.

This reporting format also makes the system practical for real use. A researcher can read the final answer quickly,
inspect the verified claims when needed, and use the trace to understand how the system reached its position.
That combination of brevity and accountability is especially valuable in an academic setting.

\section{Graceful Degradation and the Optional LLM Layer }

A defining architectural property of AI-Scientist is that the language model is an optional accelerator rather than
a hard dependency. The Verification and Synthesis agents consult a provider-backed model only when one is
configured; otherwise the system runs entirely on its deterministic heuristics. The model layer is also written to
tolerate provider differences: reasoning-capable models receive an effort parameter, while chat models that reject
that parameter are detected and called without it, and any model failure is logged and falls back to the
deterministic result.

This design serves three purposes. First, it makes the core pipeline reproducible and offline-capable, which is
essential for a thesis that emphasizes inspectable reasoning rather than opaque cloud orchestration. Second, it
lets the model improve fluency and borderline verdicts without ever being able to override the grounding and
honesty rules, because those rules run first and constrain what the model sees. Third, it keeps the architecture
model-agnostic, so the system is not tied to a single vendor.

In practice, this makes the system attractive for both demonstration and research. The deterministic core provides
reliability, while the optional language-model layer adds fluency and polish where appropriate. The result is a
balanced design that remains stable even when external services are unavailable.

\section{Why the Architecture Matches the Thesis Claim }

The AI-Scientist architecture is a direct structural expression of the thesis claim rather than a simple sequence
of implementation modules. The system begins by grounding reasoning in layered, source-prioritized evidence
rather than trusting the model's parametric memory, decomposes that evidence into discrete claims, verifies each
claim against evidence with explicit confidence, criticizes the weak claims and re-verifies them adaptively, and
records a trace of the whole process. In doing so, it embodies four commitments:

\begin{enumerate}
\item \textbf{Evidence grounding:} reasoning is bound to retrieved sources, not to generation.
\item \textbf{Calibrated confidence:} every verdict carries a transparent confidence derived from evidence
strength, corroboration, and independence.
\item \textbf{Self-critique:} a separate Critic agent flags weak claims and can demand more evidence.
\item \textbf{Traceability:} the orchestration trace and per-claim evidence make every conclusion auditable.
\end{enumerate}

These commitments distinguish AI-Scientist from both single-pass readers and standard RAG pipelines while keeping
the design compact and academically coherent.

They also explain why the architecture is well matched to the thesis goal. Each major design choice maps to a
research concern: grounding, confidence calibration, critique, or traceability. That alignment makes the chapter a
natural bridge between the problem definition and the implementation details that follow.

\section{Chapter Conclusion }

This chapter described the AI-Scientist system design as a structured, orchestrated verification pipeline rather
than a generic language-model application. The architecture begins with layered retrieval, moves through claim
extraction and deterministic evidence-bound verification, narrows uncertainty through the Critic agent and a
targeted re-verification loop, and ends with synthesis and a traceable report. The language model is an optional
refinement layer that never overrides the grounding and honesty rules. With the architecture now established, the
next chapter maps these agents to the actual implementation.

